\section{Divisions on a cyclic cover}
\label{geom:section}

The odd-degree splitting procedure will repeatedly divide divisor classes by
$1-\sigma$, where $\sigma$ is a cyclic automorphism of a curve.
Its first requirement is an existence statement:
all the needed divisions must be rational over an extension of polynomial
degree.  Its second requirement is a test that eventually distinguishes two
ramification points.  We establish the geometric inputs here.
The essential device is
to retain a divisor supported at infinity together with a chosen division
of its class in the Jacobian.  Equality of classes may lose information about that divisor;
retaining the divisor makes a final divisibility obstruction visible.

\subsection{The curve and its infinity lattice}
\label{geom:setup}

Let $K$ be a finite field of characteristic $p$, let $q\ne p$ be an odd
prime, and suppose that $K$ contains a primitive $q$-th root of unity $\zeta$.
Let $F\in K[X]$ be monic and square-free, of degree $N\ge q$ divisible by $q$,
and suppose that all its roots $x_1,\ldots,x_N$ belong to $K$.
Consider the smooth projective curve $C/K$ with function field
\[
                  K(C)=K(X,Y),\qquad Y^q=F(X).
\]
The roots $x_i$ serve only as mathematical indices here; their values need
not be known to an algorithm.

A simple root of $F$ has valuation one, so $F$ is not a $q$-th power even
over $\overline K(X)$.  Thus $C$ is geometrically integral and the map
$X:C\longrightarrow\mathbb P^1$ has degree $q$.  The affine equation is
smooth: a singularity would have $Y=0$ and $F(X)=F'(X)=0$.
At infinity, put
\[
 z=X^{-1},\qquad U=Yz^{N/q},\qquad
 U^q=z^N F(z^{-1}).
\]
At $z=0$ the right side is $1$, and its $q$ roots are simple.
There are therefore exactly $q$ rational points at infinity,
\[
 I=\{\infty_j:0\le j<q\},\qquad U(\infty_j)=\zeta^j.
\]
The ramification points are $P_i=(x_i,0)$.  Each is totally and tamely
ramified, and there is no other ramification.  The tame Riemann--Hurwitz
formula~\cite[Tags 0C1B and 0C1F]{StacksProject} consequently gives
\begin{equation}
                  g(C)=\frac{(q-1)(N-2)}2.
\label{geom:genus}
\end{equation}

Let $\sigma$ send $(X,Y)$ to $(X,\zeta Y)$ on points.
We let it act on functions by
\begin{equation}
                 (\sigma h)(X,Y)=h(X,\zeta^{-1}Y),
\label{geom:function-action}
\end{equation}
so that $\operatorname{div}(\sigma h)=\sigma\operatorname{div}(h)$.
In particular, $\sigma P_i=P_i$ and
$\sigma\infty_j=\infty_{j+1}$, with subscripts taken modulo $q$.

Write $J=\operatorname{Pic}^0(C_{\overline K})$ for the group of
geometric degree-zero divisor classes, and put
\[
 \Lambda=\left\{\sum_{j=0}^{q-1}a_j\infty_j:
                  a_j\in\mathbb Z,\ \sum_j a_j=0\right\}.
\]
For a degree-zero divisor $D$, its class is denoted by $[D]$.
The orbit sum $1+\sigma+\cdots+\sigma^{q-1}$ acts as zero on both $J$
and $\Lambda$.  For $J$, this follows because the orbit sum of a divisor
is the pullback of its pushforward to $\mathbb P^1$, and a degree-zero
divisor on $\mathbb P^1$ is principal.  For $\Lambda$, it follows by
summing its coefficients.

Let $R=\mathbb Z[\xi]$, where $\xi$ is an abstract primitive $q$-th root
of unity, and let $\xi$ act as $\sigma$.  The preceding orbit-sum identity
makes $J$ and $\Lambda$ into $R$-modules.  Set
\begin{equation}
             \lambda=1-\xi,\qquad e=(\sigma-1)\infty_0.
\label{geom:lambda}
\end{equation}
The differences $\sigma^j e$ generate $\Lambda$.  The resulting map
$R\longrightarrow\Lambda$, $a\longmapsto ae$, is an isomorphism:
it is surjective between free abelian groups of the same rank $q-1$.
In particular, multiplication by $\lambda$ is injective on $\Lambda$.
On divisors and divisor classes, multiplication by $\lambda$ means
$1-\sigma$.

We shall use
\begin{equation}
          q=u\lambda^{q-1}\quad\text{for some }u\in R^*,
          \qquad R/(\lambda)=\mathbb F_q.
\label{geom:cyclotomic}
\end{equation}
Indeed, $q=\prod_{a=1}^{q-1}(1-\xi^a)$, and each quotient
$(1-\xi^a)/(1-\xi)$ is an algebraic integer of norm one, hence a unit.
The quotient-ring assertion follows by setting $\xi=1$ in the
cyclotomic polynomial $1+T+\cdots+T^{q-1}$.
Multiplication by $\lambda$ is surjective on $J$.
To see this, identify $J$ with the geometric points of the Jacobian of
$C$~\cite[Theorem 1.1]{MilneJV86}.  Multiplication by $q$, which is prime
to the characteristic, is a surjective
isogeny~\cite[Theorem 8.2]{MilneAV86}.
Equation~\eqref{geom:cyclotomic} then implies the
surjectivity of multiplication by $\lambda$.

\subsection{Recording a class together with an infinity divisor}

Ramification divisors and the deck-transformation operator $1-\sigma$
also underlie explicit
descent on Jacobians of cyclic covers; see
Poonen and Schaefer~\cite[Section 6, Proposition 6.2]{PoonenSchaefer97}.
Here we record the infinity divisor as an additional coordinate, so that
successive divisions retain its integral divisibility information.

For $t\ge1$, define the $R$-module
\begin{equation}
 T_t=
 \frac{\{(d,W)\in J\times\Lambda:\lambda^t d=[W]\}}
      {\{([W'],\lambda^t W'):W'\in\Lambda\}}.
\label{geom:T-definition}
\end{equation}
A pair in the numerator will also denote its image in the quotient when
there is no ambiguity.  The quotient identifies changes obtained by adding
an infinity divisor to a representative of $d$.  Keeping $W$ itself,
instead of only its class, will allow us to test divisibility in the free
module $\Lambda$.

\begin{lemma}
\label{geom:filtration}
The module $T_t$ is killed by $\lambda^t$.  The map
\[
             T_t\longrightarrow T_{t+1},\qquad
             (d,W)\longmapsto(d,\lambda W)
\]
is injective, and its image is $T_{t+1}[\lambda^t]$.
Consequently, inside any $T_m$ the successive modules are precisely
$T_j=T_m[\lambda^j]$ for $1\le j\le m$.
\end{lemma}

\begin{proof}
For a qualifying pair, multiplication by $\lambda^t$ gives
$([W],\lambda^t W)$, which is one of the quotient relations.
The displayed map respects both the qualifying equation and the relations.
If its value is zero, then
\[
             (d,\lambda W)=([W'],\lambda^{t+1}W')
\]
for some $W'\in\Lambda$.  Injectivity of $\lambda$ on the lattice gives
$W=\lambda^tW'$, so the original pair was already zero in $T_t$.
Conversely, suppose that the class of $(d,W)$ in $T_{t+1}$ is killed by
$\lambda^t$.  There is then a $W'\in\Lambda$ such that
\[
       \lambda^t d=[W'],\qquad
       \lambda^t W=\lambda^{t+1}W'.
\]
The second equation gives $W=\lambda W'$, and $(d,W)$ is the image of
$(d,W')\in T_t$.  Iteration proves the last assertion.
\end{proof}

The first module has an explicit description in terms of the ramification
points.  This description will serve twice: it controls Frobenius on all
higher modules, and it provides the labels used to distinguish roots.

\begin{lemma}[Ramification generators and labels]
\label{geom:ramification}
In the setup of Section~\ref{geom:setup}, the elements
\[
                   \epsilon_i=([P_i-\infty_0],e)
                   \quad(1\le i\le N)
\]
generate $T_1$ over $\mathbb F_q$, and
\begin{equation}
                         \sum_{i=1}^N\epsilon_i=0.
\label{geom:sum-epsilon}
\end{equation}
More explicitly, let $D$ be a geometric degree-zero divisor, let
$W\in\Lambda$, and suppose that
\[
                  \operatorname{div}(h)=(1-\sigma)D-W.
\]
Then $\operatorname{Norm}_{\overline K(C)/\overline K(X)}(h)$ is a
nonzero constant.  Choose $\gamma\in\overline K^*$ with
\[
 \gamma^q=\operatorname{Norm}_{\overline K(C)/\overline K(X)}(h).
\]
Then $h(P_i)$ is defined and nonzero, and the residues
\begin{equation}
 l_i=v_{P_i}(D)-\log_\zeta\bigl(h(P_i)/\gamma\bigr)
                   \quad\text{in }\mathbb F_q
\label{geom:labels}
\end{equation}
satisfy
\begin{equation}
                         ([D],W)=\sum_{i=1}^N l_i\epsilon_i
                         \quad\text{in }T_1.
\label{geom:label-expansion}
\end{equation}
Here $v_{P_i}(D)$ is the coefficient of $P_i$ in $D$, and
$\log_\zeta$ is the unique exponent modulo $q$ of a $q$-th root of unity.
\end{lemma}

\begin{proof}
Since $\lambda[P_i-\infty_0]=[e]$, each $\epsilon_i$ qualifies.
The divisor of $Y$ is
\[
                 \operatorname{div}(Y)=\sum_iP_i-
                                      \frac Nq\sum_j\infty_j.
\]
Thus, with
\[
                   W'=\frac Nq\sum_j\infty_j-N\infty_0,
\]
we have $\sum_i[P_i-\infty_0]=[W']$ and $\lambda W'=Ne$.
This is exactly the relation proving~\eqref{geom:sum-epsilon}.

Now take any pair $([D],W)$ in the numerator defining $T_1$, and choose
$h$ with the stated divisor.  The orbit sum of $(1-\sigma)D-W$ is zero,
so the norm of $h$ has zero divisor on $\mathbb P^1$ and is a nonzero
constant.  Hilbert's Theorem~90~\cite[Example 2.3.4]{GilleSzamuely06}
for the cyclic extension
$\overline K(C)/\overline K(X)$ gives
\begin{equation}
                          h/\gamma=b/(\sigma b)
\label{geom:hilbert90}
\end{equation}
for some $b\in\overline K(C)^*$.
Choose an integral divisor $W_0$ supported on $I$ such that
$(1-\sigma)W_0=W$.  Such a divisor exists because the image of
$1-\sigma$ on the full permutation lattice $\mathbb Z^I$ is its
augmentation lattice $\Lambda$.  It need not have degree zero.
Taking divisors in~\eqref{geom:hilbert90} shows that
\[
                      D-\operatorname{div}(b)-W_0
\]
is invariant under $\sigma$.

An invariant divisor has equal coefficients along each unramified orbit.
At a ramification point its coefficient may be reduced modulo $q$, because
pullback multiplies that coefficient by $q$.  Hence
\begin{equation}
           D-\operatorname{div}(b)-W_0
                =\sum_i a_iP_i+X^*D_0,
          \qquad 0\le a_i<q,
\label{geom:invariant-divisor}
\end{equation}
for a divisor $D_0$ on $\mathbb P^1_{\overline K}$.
Put
\[
       W'=W_0+\Bigl(\sum_i a_i\Bigr)\infty_0
                       +(\deg D_0)\sum_j\infty_j.
\]
Taking degrees in~\eqref{geom:invariant-divisor} gives $\deg W'=0$.
A divisor of degree $\deg D_0$ on $\mathbb P^1$ is linearly equivalent
to $(\deg D_0)\infty$, so
\[
       [D]=[W']+\sum_i a_i[P_i-\infty_0],\qquad
       \lambda W'=W-\sum_i a_i e.
\]
The pair $([W'],\lambda W')$ vanishes in $T_1$, proving generation
and the expansion with coefficients $a_i$ modulo $q$.

It remains to identify these coefficients from $D$ and $h$.
At $P_i$, both $(1-\sigma)D$ and $W$ have coefficient zero.
Thus $h$ is a unit there.  Since $\sigma$ fixes $P_i$, evaluation of the
norm gives
\[
                          h(P_i)^q=\operatorname{Norm}(h).
\]
In particular $h(P_i)/\gamma$ is a power of $\zeta$.
The function $Y$ is a local parameter at $P_i$.  If
$v_{P_i}(b)=a$, the action~\eqref{geom:function-action} gives
\[
                         \bigl(b/(\sigma b)\bigr)(P_i)=\zeta^a.
\]
On the other hand, taking the coefficient of $P_i$
in~\eqref{geom:invariant-divisor} gives
$a_i\equiv v_{P_i}(D)-a\pmod q$.
Equation~\eqref{geom:hilbert90} now yields~\eqref{geom:labels} and
\eqref{geom:label-expansion}.
\end{proof}

In particular, if all the labels in~\eqref{geom:labels} agree, then
$([D],W)=0$ in $T_1$.  This implication does not require the labels to
be unique.  Replacing $\gamma$ by another $q$-th root shifts all of them
by the same constant, which has no effect because of
\eqref{geom:sum-epsilon}.

\subsection{Frobenius and rational divisions}

Let $\Pi$ denote the $|K|$-power Frobenius acting on geometric points,
divisors, and divisor classes.  It commutes with $\sigma$ and fixes
$\Lambda$ pointwise.  Lemma~\ref{geom:ramification} shows that it also
fixes $T_1$ pointwise, since all $P_i$ and all infinity points are
$K$-rational.  We next propagate this information through the filtration.

\begin{lemma}
\label{geom:frobenius}
For the curve and modules above, let $r\ge0$, set
$m=1+(q-1)r$, and let $k/K$ be the extension of degree $q^r$.
Then $\Pi^{q^r}$ fixes $T_m$ pointwise.  If $1\le t\le m$ and
$d\in J$, $W\in\Lambda$ satisfy $\lambda^t d=[W]$, then $d$ itself
is fixed by the Frobenius of $k$.
\end{lemma}

\begin{proof}
On $T_m$, let $F_j=\ker(\lambda^j)$ for $0\le j\le m$, and put
$F_j=0$ for negative $j$.  By Lemma~\ref{geom:filtration}, $F_1=T_1$.
For $1\le j\le m$ and $v\in F_j$,
\[
           \lambda^{j-1}(\Pi-1)v
                       =(\Pi-1)\lambda^{j-1}v=0.
\]
Thus $\Pi-1$ lowers the filtration by one.
Multiplication by $q$ lowers it by $q-1$, by
\eqref{geom:cyclotomic}.

Suppose that an endomorphism $A=\Pi^{q^a}-1$ lowers the filtration by
at least $f\ge1$.  The binomial expansion gives
\[
          \Pi^{q^{a+1}}-1
             =qA+\binom q2 A^2+\cdots+
                       \binom q{q-1}A^{q-1}+A^q.
\]
Every intermediate coefficient is divisible by $q$.
Its term therefore lowers by at least $f+q-1$; the last term lowers by
$qf\ge f+q-1$.  Induction shows that $\Pi^{q^r}-1$ lowers by
$1+r(q-1)=m$ and hence vanishes on $T_m$.

For the final assertion, $T_t$ embeds in $T_m$, so the class of $(d,W)$
in $T_t$ is fixed.  The lattice component is unchanged.  Consequently,
for some $W'\in\Lambda$,
\[
             (\Pi^{q^r}d-d,0)=([W'],\lambda^tW').
\]
The second coordinate and injectivity of $\lambda$ on $\Lambda$ force
$W'=0$.  The first coordinate then says $\Pi^{q^r}d=d$.
\end{proof}

For computation we need divisors and functions over $k$, not just fixed
geometric classes.  The following finite-field descent fact supplies them.
We include its proof to make the distinction explicit.

\begin{lemma}[Finite-field descent]
\label{geom:descent}
Let $C$ be a smooth projective geometrically integral curve over a finite
field $k$.  A $k$-rational divisor that is geometrically principal is the
divisor of a function in $k(C)^*$.  Every geometric divisor class fixed by
the Frobenius of $k$ has a $k$-rational divisor representative.
\end{lemma}

\begin{proof}
For the first assertion, choose a function $h$ with the given divisor over
a finite constant extension $k'/k$.  If $\tau$ generates
$\operatorname{Gal}(k'/k)$, then $\tau h/h$ is a constant in $(k')^*$,
and its norm to $k$ is one.  Hilbert's Theorem~90 for $k'/k$ lets us
scale $h$ by a constant so that it is fixed by $\tau$.

For the second assertion, choose a representative $D$ over some finite
extension $k'/k$.  Since its class is fixed,
$D-\tau D$ is principal.  The first assertion, applied over $k'$,
supplies $h\in k'(C)^*$ with
\[
                         D-\tau D=\operatorname{div}(h).
\]
The norm of $h$ in the constant-field extension has zero divisor, hence
is a constant in $k^*$.  The norm $(k')^*\to k^*$ is surjective, so a
constant multiple of $h$ has norm one.  Hilbert's Theorem~90 for
$k'(C)/k(C)$ now gives $h=b/(\tau b)$ after that scaling.
It follows that $D-\operatorname{div}(b)$ is invariant under $\tau$.
An invariant divisor descends by grouping its points into Galois orbits,
and its class is the class of $D$.
\end{proof}

\begin{corollary}
\label{geom:lifts}
In the setup of Section~\ref{geom:setup}, let $r\ge0$,
$m=1+(q-1)r$, and $[k:K]=q^r$.
For every $i$ there exist classes $d_{i,1},\ldots,d_{i,m}$, all represented
by degree-zero divisors over $k$, such that
\begin{equation}
 \begin{split}
 d_{i,1}&=[P_i-\infty_0],\\
 \lambda d_{i,t+1}&=d_{i,t}\quad(1\le t<m),\\
 \lambda^t d_{i,t}&=[e]\quad(1\le t\le m).
 \end{split}
\label{geom:lift-chain}
\end{equation}
Moreover, every successive choice of a geometric division in such a chain
has a representative over $k$.
\end{corollary}

\begin{proof}
Surjectivity of $\lambda$ on $J$ permits each successive choice over
$\overline K$.  The identity $\lambda d_{i,1}=[e]$ and induction give
the last equation in~\eqref{geom:lift-chain}.
Thus $(d_{i,t},e)$ qualifies for $T_t$.
Lemma~\ref{geom:frobenius} fixes $d_{i,t}$ over $k$, and
Lemma~\ref{geom:descent} supplies its divisor representative over $k$.
The argument applies to every chosen division, not merely to a special
compatible family.
\end{proof}

\subsection{Why the labels must separate ramification points}

Set $r=v_q(N)$ in Corollary~\ref{geom:lifts} and sum the last members of
its lift families.  The resulting class $d_m=\sum_i d_{i,m}$ satisfies
$\lambda^m d_m=[Ne]$.  This single relation is enough to force a
distinction between ramification points; the particular choices of lifts
no longer matter.  We now prove that assertion before turning to the
algorithms that construct the divisors and evaluate their labels.

\begin{proposition}[Forced separation]
\label{geom:separation}
In the setup of Section~\ref{geom:setup}, put $r=v_q(N)$ and
$m=1+(q-1)r$, and let $k/K$ have degree $q^r$.
Let $d_m\in J$ have a degree-zero divisor representative over $k$ and
satisfy $\lambda^m d_m=[Ne]$.
For $1\le t\le m$, set
\begin{equation}
           d_t=\lambda^{m-t}d_m,\qquad
           W_t=\lambda^{-(t-1)}Ne\in\Lambda,
\label{geom:separation-data}
\end{equation}
and choose any degree-zero divisor $D_t$ over $k$ representing $d_t$.
The lattice divisions defining $W_t$ exist and are unique.

Starting at $t=1$, make the following test, advancing to $t+1$ only when
all labels at $t$ agree.  Choose $h_t\in k(C)^*$ and $\gamma_t\in k^*$
such that
\begin{equation}
 \operatorname{div}(h_t)=(1-\sigma)D_t-W_t,
       \qquad \gamma_t^q=\operatorname{Norm}(h_t),
\label{geom:separation-functions}
\end{equation}
and form the labels
\begin{equation}
            l_{i,t}=v_{P_i}(D_t)
               -\log_\zeta\bigl(h_t(P_i)/\gamma_t\bigr)
                    \quad\text{in }\mathbb F_q.
\label{geom:separation-labels}
\end{equation}
At every reached test these choices exist over $k$ and the labels are
defined.  For some reached $t\le m$, the labels are not all equal.
This conclusion holds for every choice of the representatives $D_t$ and
of the functions and roots in~\eqref{geom:separation-functions}.
\end{proposition}

\begin{proof}
Write $N=q^r a$ with $q\nmid a$.  By~\eqref{geom:cyclotomic},
$q^r$ is an associate of $\lambda^{m-1}$, whereas the image of $a$ in
$R/(\lambda)=\mathbb F_q$ is nonzero.  Thus
\begin{equation}
       Ne\in\lambda^{m-1}\Lambda\setminus\lambda^m\Lambda.
\label{geom:lattice-obstruction}
\end{equation}
This proves that all the $W_t$ in~\eqref{geom:separation-data} exist
uniquely, and that $W_m\notin\lambda\Lambda$.

The relation $\lambda d_t=[W_t]$ is not asserted at every level in
advance.  It is the invariant that permits a test to be reached.
For $t=1$, it follows from $\lambda^m d_m=[Ne]$.
Suppose that test $t$ is reached with this invariant.
Then $(1-\sigma)D_t-W_t$ is a principal divisor over $k$.
Lemma~\ref{geom:descent} supplies $h_t\in k(C)^*$ with this divisor.
Lemma~\ref{geom:ramification} shows that $h_t$ is a unit at every $P_i$
and that
\[
                   \operatorname{Norm}(h_t)=h_t(P_i)^q\in k^*.
\]
Hence a root $\gamma_t$ exists in $k$, and every ratio
$h_t(P_i)/\gamma_t$ is a $q$-th root of unity.  The labels are defined
and satisfy
\[
                       (d_t,W_t)=\sum_i l_{i,t}\epsilon_i
                       \quad\text{in }T_1.
\]

If these labels are not all equal, the conclusion holds.
If they agree, Equation~\eqref{geom:sum-epsilon} makes this pair zero.
The quotient defining $T_1$ therefore gives an actual infinity divisor
$W'\in\Lambda$ with
\begin{equation}
                         d_t=[W'],\qquad W_t=\lambda W'.
\label{geom:separation-zero-pair}
\end{equation}
For $t<m$, injectivity of $\lambda$ on $\Lambda$ identifies $W'$ with
$W_{t+1}$.  Since $\lambda d_{t+1}=d_t$, we obtain
$\lambda d_{t+1}=[W_{t+1}]$, exactly the invariant for the next test.
At $t=m$, however, \eqref{geom:separation-zero-pair} contradicts
$W_m\notin\lambda\Lambda$.  Not all tests can have equal labels.
The argument used no restrictions on the representatives or the choices
in~\eqref{geom:separation-functions}, proving the final assertion.
\end{proof}

Figure~\ref{geom:ladder} separates the two traversals in this proof.
The class $d_m$ determines the class chain by applying $\lambda$; the
divisor $Ne$ determines the infinity chain by dividing by $\lambda$.
The label tests compare their progress and would force one impossible
extra lattice division if every test were constant.

\begin{figure}[htbp]
\centering
\begin{tikzpicture}[>=Stealth, every node/.style={font=\small},
                    point/.style={minimum width=1.6cm,minimum height=0.7cm}]
 \node at (-1.5,0) {$J$};
 \node at (-1.5,-1.5) {$\Lambda$};
 \node[point] (d1) at (0,0) {$d_1$};
 \node[point] (d2) at (2.6,0) {$d_2$};
 \node[point] (dd) at (5.2,0) {$\cdots$};
 \node[point] (dm) at (7.8,0) {$d_m$};
 \draw[->] (d2) -- node[above] {$\lambda$} (d1);
 \draw[->] (dd) -- node[above] {$\lambda$} (d2);
 \draw[->] (dm) -- node[above] {$\lambda$} (dd);
 \node[point] (w1) at (0,-1.5) {$W_1=Ne$};
 \node[point] (w2) at (2.6,-1.5) {$W_2$};
 \node[point] (ww) at (5.2,-1.5) {$\cdots$};
 \node[point] (wm) at (7.8,-1.5) {$W_m$};
 \draw[->] (w1) -- node[above] {divide by $\lambda$} (w2);
 \draw[->] (w2) -- node[above] {divide by $\lambda$} (ww);
 \draw[->] (ww) -- node[above] {divide by $\lambda$} (wm);
 \node[below=1mm of wm] {$W_m\notin\lambda\Lambda$};
\end{tikzpicture}
\caption{The two chains in Proposition~\ref{geom:separation}.
Tests run from left to right.  Reaching test $t$ requires
$\lambda d_t=[W_t]$; equal labels at $t<m$ give
$d_t=[W_{t+1}]$ and allow the next test.  Equal labels at the final test
would require the forbidden additional division of $W_m$ in $\Lambda$.
No equality between the rows is assumed before its test is reached.}
\label{geom:ladder}
\end{figure}

The geometric mechanism is now complete: Corollary~\ref{geom:lifts}
supplies the class $d_m$, and Proposition~\ref{geom:separation} forces
distinct labels using that class.  The next sections implement divisor
arithmetic, class division, and evaluation without recovering the roots
of $F$ or factoring the finite supports of divisors.
